\subsection{Further obstructions from proper quotients}

\begin{proposition}[No good-outside-three abelian quotient]
\label{prop:lm-no-good-outside-three-quotient}
If $3\nmid n$, the Jacobian of $X_n$ has no positive-dimensional
abelian quotient over $\mathbb Q$ with good reduction outside
$3$.  In particular, there is no nonconstant morphism over
$\mathbb Q$ from $X_n$ to a genus-one curve whose Jacobian has
good reduction outside $3$.
\end{proposition}
\begin{proof}
The Fermat curve and its Jacobian have good reduction at $3$.
Good reduction passes to abelian quotients by the
N\'eron--Ogg--Shafarevich criterion.  A quotient as in the
statement would therefore have good reduction everywhere
over $\mathbb Q$, contradicting Fontaine's theorem that no
positive-dimensional abelian variety over $\mathbb Q$ has
everywhere good reduction.  A nonconstant map to a
genus-one curve induces a surjection onto its Jacobian.
\end{proof}

For this input see \citet[Corollary, p.~517]{Fontaine1985}.

\subsection{A local obstruction to lifting rational points of the three-punctured line}
\label{subsec:lm-local-metabelian-separation}

\begin{proposition}[Local point/cusp separation]
\label{prop:lm-local-metabelian-separation}
Let $U_1=\mathbb P^1\setminus\{0,1,\infty\}$ over
$\mathbb Q_3$.  No ordinary point section of $U_1(\mathbb Q_3)$
is cuspidal in the geometrically pro-$3$ metabelian quotient.
The three rational cusp packets in that quotient are pairwise
disjoint.
\end{proposition}
\begin{proof}
Fix $a\in\mathbb Q_3\setminus\{0,1\}$ and consider
\[
 V_a:\quad a x^9+(1-a)y^9=1,\qquad xy\ne0,
 \qquad t=a x^9.
\]
This is a finite \'etale cover of $U_1$, geometrically Galois
with deck group $(\mathbb Z/9\mathbb Z)^2$.  It has the point
$Q=(1,1)$ above $a$, and its compactification $\overline V_a$
is geometrically a Fermat curve of degree $9$.

Put $N=\pi_1(U_{1,\overline{\mathbb Q}_3})^{\mathrm{pro}\text{-}3}$.
The geometric group $H$ of $V_a$ embeds as an open subgroup of
$N$.  Since $N/H$ is abelian, the closed derived groups satisfy
\[
 N'\subset H,\qquad N''\subset H'.
\]
Let $\mathcal H$ be the arithmetic subgroup corresponding to
$V_a$, using $Q$ to lift the section of $a$.

Suppose that the section of $a$ is cuspidal modulo $N''$.
After lifting its conjugator to $N$, its full point section
and a conjugate of a full cusp section differ elementwise by
$N''$.  Because $N''\subset\mathcal H$, the conjugated cusp
section is itself contained in $\mathcal H$.  The intersection
of its cusp decomposition group with $\mathcal H$ is a cusp
decomposition group of the cover and maps onto
$G_{\mathbb Q_3}$.  Hence its boundary center
$P\in\overline V_a\setminus V_a$ is rational over
$\mathbb Q_3$.

The two sections agree after passing from $H$ to $H^{\mathrm{ab}}$,
since $N''\subset H'$.  After compactification their Kummer
classes in $H^1(\mathbb Q_3,T_3\operatorname{Jac}(\overline V_a))$
agree.  The local completion kernel calculation in the proof of Lemma~\ref{lem:fourzero} shows that
$[Q-P]$ is torsion of order prime to $3$.  \citet[Theorem~A, p.~191]{Coleman1986} applies to the geometric degree-nine Fermat curve,
since $9=(19-1)/2$ and $2\le8$.  It forces a point
torsion-equivalent to a coordinate cusp to be a coordinate
cusp.  This contradicts the fact that $Q=(1,1)$ is interior.

For the second assertion, use the two Kummer coordinates
$t$ and $1-t$ to identify the geometric abelianization with
$\mathbb Z_3(1)^2$.  The three cusp packets there are the
lines with directions $(1,0)$, $(0,1)$ and $(-1,-1)$.
Here constants $-1$ have trivial pro-$3$ Kummer class.
Two different packets intersect only in the zero class.
If two full cusp sections agreed modulo $N''$, both cusp
parameters would therefore be zero.  They lift to the
standard cover $x^3+y^3=1$.  As above, equality modulo
$N''$ persists in the geometric abelianization of this
cover, because its geometric deck group is abelian.
Their compact images are two different rational coordinate
cusps of the smooth cubic.  Their difference is a nonzero
point of order $3$ on its elliptic Jacobian.  Such a point
is not in the prime-to-$3$ torsion kernel of local completion,
a contradiction.
\end{proof}

\begin{corollary}
\label{cor:lm-no-rational-point-lift}
Suppose that every geometrically pro-$3$ Selmer section of
$U_n/\mathbb Q$ is locally cuspidal at $3$, as in
Theorem~\ref{thm:localthreeintro}.  Its pushdown under the
coordinate map $U_n\to U_1$ cannot be the section of an
ordinary rational point of $U_1$, even if this equality is
requested only in the geometrically pro-$3$ metabelian quotient.
In particular, the rational point $t=2$ gives no such Selmer lift.
\end{corollary}
\begin{proof}
Localizing a putative equality at $3$ would make the ordinary
point section cuspidal, contrary to
Proposition~\ref{prop:lm-local-metabelian-separation}.
\end{proof}
